% =====================================================================
%  ThmSmith Stage -1 proposal skeleton.
%  Written automatically by the ThmSmith TS pipeline before Stage -1 runs.
%  Locked parameters: qid=eid\_proximal\_phase, specialization=v1
%
%  HOW TO USE THIS FILE
%  --------------------
%  - The preamble (everything above the `% --- BEGIN CONTENT ---` marker)
%    and the `\section{...}` headers below are fixed by the pipeline.
%    Do NOT rewrite or rename them; the Step 3 reviewer subagent and the
%    downstream Stage 0 / Stage 1 prompts grep for these section titles.
%  - Each section contains exactly one `% TODO(stage-1 ...): ...`
%    placeholder. Replace each TODO with the content described for that
%    section in `ThmSmith/tools/src/prompts/stage_neg1_draft.txt`
%    ("REQUIRED OUTPUT FORMAT (Stage -1 proposal .tex)").
%  - The `\paragraph{Locked parameters.}` block in the metadata block
%    must pin the chosen `(qid, specialization, p, assignment, target,
%    regression)` precisely. If the proposal maps to the Q1 pool-mode,
%    reproduce that block verbatim from
%    `ThmSmith/doc/panel_formula_question_generator_note_with_common_prompts.tex`
%    (Q2..Q6 templates have been retired). Otherwise (free-form
%    `(qid, specialization)` — the default), define each parameter
%    explicitly here (no implicit conventions). Stage 0 quotes this block;
%    do not rephrase it after locking.
%  - Removing a section header, renaming a macro, or rewriting the
%    preamble is a regression: the file must still match this skeleton
%    after you finish.
% =====================================================================

\documentclass[11pt]{article}

\usepackage{amsmath,amssymb,amsthm,mathtools}
\usepackage{enumitem}
\usepackage[colorlinks=true,linkcolor=blue,citecolor=blue,urlcolor=blue]{hyperref}
\usepackage{natbib}

\newcommand{\E}{\mathbb{E}}
\renewcommand{\Pr}{\mathbb{P}}
\newcommand{\one}{\mathbf{1}}
\newcommand{\transpose}{\mathsf{T}}
\newcommand{\calH}{\mathcal{H}}
\newcommand{\calF}{\mathcal{F}}
\newcommand{\calR}{\mathcal{R}}
\newcommand{\R}{\mathbb{R}}
\newcommand{\plim}{\operatorname*{plim}}

\theoremstyle{definition}
\newtheorem{definition}{Definition}
\newtheorem{assumption}{Assumption}
\newtheorem{example}{Example}

\theoremstyle{plain}
\newtheorem{theorem}{Theorem}
\newtheorem{conjecture}{Conjecture}
\newtheorem{proposition}{Proposition}
\newtheorem{lemma}{Lemma}
\newtheorem{corollary}{Corollary}

\theoremstyle{remark}
\newtheorem{remark}{Remark}

\title{ThmSmith Stage -1 proposal: \texttt{eid\_proximal\_phase} / \texttt{v1}%
  \\ \large\textit{A native-equivalence frontier for finite proximal OR/DR and minimax bridge estimators}}
\author{ThmSmith Stage -1 proposer}
\date{\today}

\begin{document}
\maketitle

% --- BEGIN CONTENT ---  (everything above is fixed by the pipeline)

\paragraph{Locked parameters.}
This is a partial-identification / exact-functional-identification proposal
with locked tuple
\[
(\texttt{qid},\texttt{specialization}, P, \theta, \text{identifying restrictions},
\text{bounding and native selection-range functionals}).
\]
Here \(\texttt{qid}=\texttt{eid\_proximal\_phase}\) and
\(\texttt{specialization}=\texttt{v1}\).  The observable distribution \(P\)
is the finite-support joint law of \(O=(Y,A,W,Z,X)\), with
\(A\in\{0,1\}\), \(Y\in[0,1]\), finite proxy supports for \(W\) and \(Z\),
and finite covariate support for \(X\).  The target parameter is the proximal
average treatment effect \(\theta=\E\{Y(1)-Y(0)\}\).  The identifying
restrictions are consistency, positivity for both treatment arms, latent
proximal exchangeability, negative-control exclusions, existence of a bounded
outcome bridge \(h(W,A,X)\), and the finite bridge-fiber sharpness condition
stated in Assumption~\ref{ass:fiber-sharp}.  The proximal OR bridge
functional is
\[
\Theta_I^{\mathrm{OR}}(P)=\{\ell_P(h):h\in\calH_P^{\mathrm{OR}}\},\qquad
\ell_P(h)=\E_P\{h(W,1,X)-h(W,0,X)\},
\]
where \(\calH_P^{\mathrm{OR}}=\{h\in[0,1]^d:T_Ph=m_P\}\),
\((T_Ph)(z,a,x)=\E_P\{h(W,a,x)\mid Z=z,A=a,X=x\}\), and
\(m_P(z,a,x)=\E_P(Y\mid Z=z,A=a,X=x)\).  The finite-class minimax bridge
objective is represented by observable matrices \((B_P,c_P,R_P)\), where
\(B_P\) and \(c_P\) come from a finite bridge dictionary and finite critic
class, \(R_P\succ0\), and the exact-fit fiber is
\(\calH_P^{\mathrm{MM}}=\{h\in[0,1]^d:B_Ph=c_P\}\); no equality
\(\calH_P^{\mathrm{MM}}=\calH_P^{\mathrm{OR}}\) is assumed.  The native
selection functionals for an admissible quadratic regularizer \(Q\succ0\)
are
\[
\theta_{\mathrm{OR},Q}(P)=\ell_P(h_{\mathrm{OR},Q}),\qquad
h_{\mathrm{OR},Q}=\operatorname*{arg\,min}_{h\in\calH_P^{\mathrm{OR}}} h^\transpose Qh,
\]
and
\[
\theta_{\mathrm{MM},Q}(P)=\ell_P(h_{\mathrm{MM},Q}),\qquad
h_{\mathrm{MM},Q}=\operatorname*{arg\,min}_{h\in\calH_P^{\mathrm{MM}}} h^\transpose Qh.
\]
The native selection ranges are
\[
\mathcal S_{\mathrm{OR}}(P)=\{\theta_{\mathrm{OR},Q}(P):Q=Q^\transpose\succ0\},
\qquad
\mathcal S_{\mathrm{MM}}(P)=\{\theta_{\mathrm{MM},Q}(P):Q=Q^\transpose\succ0\}.
\]
Write
\(\Theta_I^j(P)=\ell_P(\calH_P^j)=[\underline\theta_j(P),\overline\theta_j(P)]\)
for \(j\in\{\mathrm{OR},\mathrm{MM}\}\).  The observable native-invariance
threshold is
\[
\tau_{\mathrm{native}}(P)=
\max\{\overline\theta_{\mathrm{OR}}-\underline\theta_{\mathrm{OR}},
\overline\theta_{\mathrm{MM}}-\underline\theta_{\mathrm{MM}},
|\underline\theta_{\mathrm{OR}}-\underline\theta_{\mathrm{MM}}|,
|\overline\theta_{\mathrm{OR}}-\overline\theta_{\mathrm{MM}}|\}.
\]
Let \(\psi_{\mathrm{OR/DR},Q,P}\) denote the finite-dimensional first-order
influence function for the proximal OR/IPW/DR plug-in estimator using
\(h_{\mathrm{OR},Q}\), and let \(\psi_{\mathrm{MM},Q,P}\) denote the
corresponding plug-in first-order influence function for the profiled minimax
bridge estimator using \(h_{\mathrm{MM},Q}\), whenever the fixed-rank
regularity conditions in Assumption~\ref{ass:regularity} hold.  The novelty
question is whether \(\tau_{\mathrm{native}}(P)=0\) is the sharp observable
boundary for these two finite native bridge objectives to agree on the same
singleton identified set, with the same first-order/efficiency-bound object
when the zero-threshold equality holds locally on a fixed-rank stratum, and
whether \(\tau_{\mathrm{native}}(P)>0\) forces regularizer dependence or
cross-workflow disagreement rather than a hidden point-identified causal value.

\begin{abstract}
Proximal causal inference often solves an ill-posed bridge equation by
regularization, even though the causal estimand is identified only if the ATE
loading is invariant over the bridge solution fiber selected by the native
population objective.  This proposal asks for the sharp population
distinction between a genuinely identified proximal ATE and values selected
by two finite native bridge objectives tied to proximal analyses when
completeness fails: the residual-squared conditional-moment bridge objective
underlying finite proximal OR/IPW/DR implementations and a profiled
finite-quadratic minimax negative-control objective.  A support
proposition records the known finite null-annihilator criterion, while the
flagship conjecture is a new observable native-objective threshold
\(\tau_{\mathrm{native}}(P)\): zero means both finite objectives are invariant and
agree on one singleton identified set, and local zero on a fixed-rank stratum
also gives the same finite first-order limit; nonzero means at least one finite objective is
penalty-dependent or the two native objectives select different causal values
under the same observed law.  A second conjecture predicts that, on an
open finite negative-control class with ATE-active bridge nonuniqueness, the
native quadratic selector ranges have nonempty interior and fill the
one-dimensional sharp interval after closure.  Resolving this would tell
applied proximal analysts when a regularized proximal point estimate is an
identification statement and when it is only a property of the chosen bridge
solver.
\end{abstract}

\section{Introduction}
\paragraph{Why the problem is important.}
Negative-control proxies let researchers use information about latent
confounding when ordinary measured-confounder adjustment is not credible.
The bridge equation is an inverse problem, and in practice researchers often
stabilize it with a penalty or sieve even when the underlying completeness
condition is doubtful.  The scientific question is therefore not only whether
a bridge can be estimated, but whether the causal value delivered by the
chosen regularizer is the same for every data-compatible bridge.

\paragraph{The main finding (proposed).}
Proposition~\ref{prop:null-annihilator} is only support: it records the
known finite population criterion for regularization-invariant proximal ATEs.
The flagship statements are Conjecture~\ref{conj:penalty-separation} and
Conjecture~\ref{conj:selection-range}.  They compare two finite native
bridge objectives without imposing a shared exact bridge fiber in advance:
the finite residual-squared bridge step in the proximal OR/IPW/DR workflow of
\citet{CuiPuShiMiaoTchetgen2024JASA} and the finite profiled-quadratic bridge
step in the minimax workflow of \citet{KallusMaoUehara2021MinimaxNC}.  The
proposed frontier is the observable LP threshold
\(\tau_{\mathrm{native}}(P)\): at zero, the two finite objectives are
invariant and agree on the singleton identified set, and if this zero-threshold
condition holds locally on the fixed-rank stratum then they have the same
first-order causal limit; above zero, either a
native selector range moves with \(Q\), the two objectives disagree on the ATE
value under the same observed law, or no common efficiency-bound
interpretation can be assigned to both selected values.

\paragraph{What is new.}
The closest published positive results, \citet{MiaoGengTchetgen2018Biometrika},
\citet{CuiPuShiMiaoTchetgen2024JASA}, and
\citet{ZhangLiMiaoTchetgen2023SPL}, show that proximal point identification
can hold under bridge/completeness or even bridge nonuniqueness when the
target functional remains invariant.  The closest weak-identification
functional result, \citet{BennettKallusMaoNeweySyrgkanisUehara2023COLT},
studies inference for strongly identified linear functionals of weakly
identified nuisance functions, and the closest minimax regularization
comparator, \citet{BennettKallusMaoNeweySyrgkanisUehara2023MinimaxIV}, proves
that a penalized minimax IV estimator can converge in \(L_2\) to a fixed
solution without uniqueness or closedness.  The epsilon gap here is the
complementary proximal-identification case: before reducing both finite
objectives to a common \(Q\)-selector, give a computable observable frontier
for whether their native bridge-solving objectives are equivalent as causal
estimators--same singleton identified set, and on a local zero-threshold
stratum the same first-order limit and efficiency-bound object in the finite
regular case--or only as
penalty-selected bridge algorithms.  The flagship content is not the
already-known
null-annihilator criterion and not the definitional fact that selected bridge
values lie inside a sharp interval.  It is the proposed iff boundary
\(\tau_{\mathrm{native}}(P)=0\) versus \(\tau_{\mathrm{native}}(P)>0\), plus
the open-class nonempty-interior result, SPD-exposed-point coverage, and
one-dimensional closure equality for the native selector ranges.

\paragraph{How it positions in the literature.}
The proposal sits at the boundary between exact proximal identification and
partial identification.  The exact-ID side supplies the bridge equation and
the null-annihilator criterion for when a nonunique bridge still identifies
the ATE.  The partial-ID side supplies the native sharp intervals
\(\Theta_I^{\mathrm{OR}}(P)\) and \(\Theta_I^{\mathrm{MM}}(P)\) once the
annihilator criterion fails for a workflow.  The inverse-problem literature
supplies residual-squared and minimax bridge selectors, but the proposal's
target is not estimation theory; it is the population identification status
of the point selected by those named regularizers.

\paragraph{Implications for the community.}
The concrete downstream consumer is the semiparametric proximal OR/IPW/DR
analysis of \citet{CuiPuShiMiaoTchetgen2024JASA} when its outcome bridge is
implemented through the finite residual-squared conditional-moment objective
specified in Assumption~\ref{ass:native-workflow}; a second consumer is the
finite-class profiled-quadratic specialization of the minimax
negative-control bridge learning problem of
\citet{KallusMaoUehara2021MinimaxNC}.  In either finite objective, when
completeness is questionable, the analyst should report whether
\(\tau_{\mathrm{native}}(P)=0\) or whether the reported proximal point
estimate depends on the native objective or penalty that selected one bridge
from many.

\paragraph{How research benefits from the conclusion.}
Researchers using proximal causal inference should check the finite
native-objective threshold \(\tau_{\mathrm{native}}(P)\) before treating a
regularized bridge estimate as a point-identified ATE.

\section{Related work}
\citet{MiaoGengTchetgen2018Biometrika} is the foundational proximal ATE
point-identification paper: it identifies causal effects with two proxy
variables under rank/completeness-type bridge conditions, leaving
regularization under failed completeness outside the theorem.
\citet{CuiPuShiMiaoTchetgen2024JASA} develops proximal OR, IPW, and doubly
robust semiparametric estimators under bridge conditions; this proposal asks
what finite bridge objectives select when several bridges solve the same observable
moment equation.  \citet{TchetgenYingCuiShiMiao2024StatSci} gives the current
overview of proximal causal inference and frames bridge assumptions as the
replacement for measured exchangeability.  \citet{ZhangLiMiaoTchetgen2023SPL}
shows that nonunique proximal bridge functions can still identify
counterfactual means; Proposition~\ref{prop:null-annihilator} records the
finite linear-algebra condition behind that already-known phenomenon, while
Conjecture~\ref{conj:penalty-separation} studies the opposite case.
\citet{BennettKallusMaoNeweySyrgkanisUehara2023COLT} studies root-\(n\)
inference on strongly identified linear functionals of weakly identified
nuisance functions, including proximal examples; the proposed conjecture
asks what regularizers select when the functional is not strongly identified.
\citet{BennettKallusMaoNeweySyrgkanisUehara2023MinimaxIV} is the closest
regularization comparator: it gives a penalized minimax IV estimator with
strong \(L_2\) convergence to a fixed solution without uniqueness or
closedness, whereas this proposal asks whether the fixed selected solution is
also the same proximal causal estimand selected by the finite OR/DR bridge
workflow.
\citet{KallusMaoUehara2021MinimaxNC} develops minimax learning of
negative-control bridge functions while deliberately avoiding strong
completeness and uniqueness requirements; this proposal does not claim to
cover their full saddle-point estimator, but targets the population causal
object selected when a finite bridge dictionary and finite critic class
reduce the profiled bridge subproblem to the quadratic matrices
\((B_P,c_P,R_P)\) in Assumption~\ref{ass:native-workflow}.
\citet{Santos2011JoE} gives the NPIV analogue separating recovery of a
linear functional from recovery of the whole structural function, and
\citet{Santos2012Econometrica} provides a partial-identification inference
anchor for conditional moment models.  \citet{GhassamiShpitserTchetgen2023ProxyPI}
develops proxy-variable partial-identification bounds without completeness;
the present proposal uses that partial-ID perspective but focuses on the
extra point selected by bridge regularization.  \citet{DarollesFanFlorensRenault2011Econometrica}
is the Econometrica anchor for Tikhonov regularization in NPIV inverse
problems; here the penalty is treated as a selector on an exact bridge fiber,
not as additional causal information.  \citet{Horowitz2014ARE} reviews
ill-posed inverse problems in econometrics and motivates why regularized
solutions can be stable estimators of penalty-dependent objects.  \citet{CanaySantosShaikh2013Econometrica}
shows that completeness and related identification hypotheses are generally
not testable in unrestricted nonparametric models, motivating the finite
observable matrix convention locked here.  \citet{Manski1990AER} supplies the
bounded-outcome partial-identification fallback that appears when the bridge
fiber can move the ATE over a nontrivial interval.

In-repository prior art does not already contain this kernel.
\texttt{ThmSmith/doc/API.md} \S3 records that the ExactID and PartialID
catalogues are empty seed areas for new exact-identification and
partial-identification qids.  The closest banked artifact is
\texttt{pid\_mediation\_proxy}, whose reviewer record warns that proxy
bridge-null geometry must not be declared sharp without a real causal
realization argument; this proposal avoids reusing that rejected sharpness
claim by making bridge-fiber sharpness an explicit assumption and studying
regularizer dependence conditional on it.  The downgraded
\texttt{eid\_frontdoor\_multi\_mediator} proposal is adjacent because it
uses proxy-envelope geometry, but its kernel is multi-mediator coupling, not
proximal ATE regularization.  The accepted
\texttt{q1\_spectral\_phase\_transition} banked analogue concerns panel/FWL
singular-value phase transitions rather than negative-control bridge
regularization.

\section{Formal setup}
Let \(O=(Y,A,W,Z,X)\) have finite observable support, with
\(A\in\{0,1\}\), \(Y\in[0,1]\), and observed law \(P\).  The latent proximal
model contains an unobserved confounder \(U\), potential outcomes \(Y(0)\)
and \(Y(1)\), consistency \(Y=Y(A)\), latent exchangeability
\((Y(0),Y(1))\perp A\mid U,X\), and negative-control exclusions saying that
\(Z\) is treatment-confounding proxy information and \(W\) is
outcome-confounding proxy information.  The finite vector space for bridges is
\[
\calH=\R^d,\qquad
d=|\mathcal W|\,|\mathcal A|\,|\mathcal X|,
\]
with coordinates \(h(w,a,x)\).  Define the observable bridge operator and
moment vector by
\[
(T_Ph)(z,a,x)=\sum_w h(w,a,x)P(w\mid z,a,x),\qquad
m_P(z,a,x)=\E_P(Y\mid Z=z,A=a,X=x).
\]
The exact bounded proximal OR bridge fiber is
\[
\calH_P^{\mathrm{OR}}=\{h\in[0,1]^d:T_Ph=m_P\}.
\]
The ATE loading is the observable linear functional
\[
\ell_P(h)=\sum_{w,x} h(w,1,x)P(w,x)-\sum_{w,x}h(w,0,x)P(w,x),
\]
and the proximal OR sharp interval is
\[
\Theta_I^{\mathrm{OR}}(P)=\{\ell_P(h):h\in\calH_P^{\mathrm{OR}}\}
=[\underline\theta_{\mathrm{OR}},\overline\theta_{\mathrm{OR}}].
\]
For a symmetric positive-definite matrix \(Q\), the native proximal OR
quadratic selector is
\[
h_{\mathrm{OR},Q}(P)=
\operatorname*{arg\,min}_{h\in\calH_P^{\mathrm{OR}}}h^\transpose Qh,\qquad
\theta_{\mathrm{OR},Q}(P)=\ell_P(h_{\mathrm{OR},Q}(P)).
\]
The proximal OR selection range is
\[
\mathcal S_{\mathrm{OR}}(P)=
\{\theta_{\mathrm{OR},Q}(P):Q=Q^\transpose\succ0\}
\subseteq\Theta_I^{\mathrm{OR}}(P).
\]
The finite-class minimax negative-control bridge objective is represented by
observable matrices \((B_P,c_P,R_P)\), with \(R_P\succ0\).  Here \(B_P\) and
\(c_P\) are the moment matrix and target vector induced by a finite bridge
dictionary and finite critic class, and \(R_P\) is the observable positive
definite critic-weight matrix after profiling the critic.  The native
profiled quadratic criterion is
\[
\min_{h\in[0,1]^d}(B_Ph-c_P)^\transpose R_P(B_Ph-c_P)+\alpha h^\transpose Qh.
\]
Its exact-fit fiber is
\[
\calH_P^{\mathrm{MM}}=\{h\in[0,1]^d:B_Ph=c_P\},
\]
which need not equal \(\calH_P^{\mathrm{OR}}\).  The minimax native selector
and range are
\[
h_{\mathrm{MM},Q}(P)=
\operatorname*{arg\,min}_{h\in\calH_P^{\mathrm{MM}}}h^\transpose Qh,\qquad
\theta_{\mathrm{MM},Q}(P)=\ell_P(h_{\mathrm{MM},Q}(P)),
\]
\[
\mathcal S_{\mathrm{MM}}(P)=
\{\theta_{\mathrm{MM},Q}(P):Q=Q^\transpose\succ0\}
\subseteq\Theta_I^{\mathrm{MM}}(P),
\quad
\Theta_I^{\mathrm{MM}}(P)=
[\underline\theta_{\mathrm{MM}},\overline\theta_{\mathrm{MM}}].
\]
The corresponding penalized finite-quadratic minimax selector is
\[
h_{\mathrm{MM},\alpha,Q}(P)=\operatorname*{arg\,min}_{h\in[0,1]^d}
(B_Ph-c_P)^\transpose R_P(B_Ph-c_P)+\alpha h^\transpose Qh,
\]
and when \(\calH_P^{\mathrm{MM}}\neq\emptyset\),
\(h_{\mathrm{MM},\alpha,Q}(P)\to h_{\mathrm{MM},Q}(P)\) as
\(\alpha\downarrow0\).
The native observable threshold is
\[
\tau_{\mathrm{native}}(P)=
\max\{\overline\theta_{\mathrm{OR}}-\underline\theta_{\mathrm{OR}},
\overline\theta_{\mathrm{MM}}-\underline\theta_{\mathrm{MM}},
|\underline\theta_{\mathrm{OR}}-\underline\theta_{\mathrm{MM}}|,
|\overline\theta_{\mathrm{OR}}-\overline\theta_{\mathrm{MM}}|\}.
\]
For \(j\in\{\mathrm{OR},\mathrm{MM}\}\), let \(F_P^j\) denote the minimal
face of \(\calH_P^j\) containing its relative interior.  The matrix \(A_P^j\)
stacks the native exact-fit equations with the active zero/one coordinates
on \(F_P^j\), and \(M_P^j\) is the Euclidean orthogonal projector onto
\(\operatorname{ker}(A_P^j)\).  The workflow-specific invariance moduli are
\[
\kappa_j(P)=\|M_P^j\ell_P^\transpose\|_2,\qquad
j\in\{\mathrm{OR},\mathrm{MM}\}.
\]
An exposed bridge point is a point \(h^\star\in\calH_P^j\) that uniquely
maximizes a linear functional over \(\calH_P^j\); its ATE value is
\(\ell_P(h^\star)\).  It is SPD-exposed for workflow \(j\) if the normal
cone of \(\calH_P^j\) at \(h^\star\) contains a vector \(-2Qh^\star\) for
some symmetric positive-definite \(Q\), the first-order condition for
\(h^\star\) to be selected by minimizing \(h^\transpose Qh\) over
\(\calH_P^j\).
The threshold is computable by four observable linear programs:
\[
\underline\theta_j(P)=\min\{\ell_P(h): h\in\calH_P^j\},\qquad
\overline\theta_j(P)=\max\{\ell_P(h): h\in\calH_P^j\},
\quad j\in\{\mathrm{OR},\mathrm{MM}\}.
\]
For a fixed admissible \(Q\), each selected value
\(\theta_{j,Q}(P)\) is computed by the strictly convex quadratic program
\(\min_{h\in\calH_P^j}h^\transpose Qh\) followed by \(\ell_P\).  For a
candidate exposed point \(h^\star\), SPD-exposure is checked by the finite
semidefinite feasibility problem \(Q=Q^\transpose\), \(Q\succeq\varepsilon I\)
for some \(\varepsilon>0\), and
\(-2Qh^\star\in N_{\calH_P^j}(h^\star)\), where the normal cone is the
polyhedral cone read from the active affine and box constraints.
The residual-penalized bridge estimator with penalty \(Q\) has population
limit
\[
h_{\mathrm{OR},\alpha,Q}(P)=\operatorname*{arg\,min}_{h\in[0,1]^d}
\|T_Ph-m_P\|_2^2+\alpha h^\transpose Qh,
\]
and when \(\calH_P^{\mathrm{OR}}\neq\emptyset\),
\(h_{\mathrm{OR},\alpha,Q}(P)\to h_{\mathrm{OR},Q}(P)\) as
\(\alpha\downarrow0\).
This exact-fiber selector is the finite-population abstraction of the
outcome-bridge step in a proximal OR analysis when the bridge is computed by
squared residuals plus an SPD ridge penalty.
More explicitly, the population residual-squared proximal OR criterion covered
here is
\[
\min_{h\in[0,1]^d}\sum_{z,a,x}\omega_{zax}
\left(m_P(z,a,x)-(T_Ph)(z,a,x)\right)^2+\alpha h^\transpose Qh,
\qquad \omega_{zax}>0,
\]
which is the displayed \(h_{\mathrm{OR},\alpha,Q}(P)\) after absorbing the positive
weights into the Euclidean norm.  No claim is made for nonsmooth,
adversarial, or nonquadratic minimax bridge limits; the minimax workflow is
included only in the finite-class profiled quadratic form above.

\begin{quote}
\begin{verbatim}
qid = eid_proximal_phase
specialization = v1
observable distribution P = finite-support law of O=(Y,A,W,Z,X)
parameter theta = E[Y(1)-Y(0)]
identifying restrictions = consistency, positivity, latent proximal exchangeability,
  negative-control exclusions, bounded outcomes, existence of a bounded outcome bridge,
  and finite bridge-fiber sharpness
OR bounding functional = Theta_I^OR(P) = {ell_P(h): h in H_P^OR}
minimax exact-fit fiber = H_P^MM = {h in [0,1]^d : B_P h = c_P}
native selection functionals =
  theta_OR,Q(P) = ell_P(argmin_{h in H_P^OR} h'Qh)
  theta_MM,Q(P) = ell_P(argmin_{h in H_P^MM} h'Qh)
native threshold = tau_native(P), the max of the two native interval widths
  and the endpoint disagreements between Theta_I^OR(P) and Theta_I^MM(P)
\end{verbatim}
\end{quote}

\section{Assumptions}
\begin{assumption}[Finite bounded observable law]\label{ass:finite-bounded}
The observable support of \((Y,A,W,Z,X)\) is finite except that
\(Y\in[0,1]\), both treatment arms have positive probability in every
covariate stratum used in \(\ell_P\), and all conditional probabilities
defining \(T_P\) are well defined.
\end{assumption}

\begin{assumption}[Proximal negative-control validity]\label{ass:proximal-valid}
There exists a latent variable \(U\) and potential outcomes \(Y(0),Y(1)\)
such that consistency, latent exchangeability given \((U,X)\), and the usual
negative-control exclusions for \(Z\) and \(W\) hold.  These restrictions are
the population restrictions that make an outcome bridge causally meaningful.
\end{assumption}

\begin{assumption}[Bounded bridge existence]\label{ass:bridge-exists}
The native bounded bridge fibers
\(\calH_P^{\mathrm{OR}}=\{h\in[0,1]^d:T_Ph=m_P\}\) and
\(\calH_P^{\mathrm{MM}}=\{h\in[0,1]^d:B_Ph=c_P\}\) are nonempty.
\end{assumption}

\begin{assumption}[Finite bridge-fiber sharpness]\label{ass:fiber-sharp}
Every \(h\) in either native bridge fiber
\(\calH_P^{\mathrm{OR}}\cup\calH_P^{\mathrm{MM}}\) is attained by at least
one finite latent proximal model satisfying
Assumptions~\ref{ass:finite-bounded}--\ref{ass:proximal-valid} and inducing
the same observed law \(P\); conversely every compatible finite latent
proximal model relevant to either native objective induces an ATE value
\(\ell_P(h)\) for some \(h\) in the corresponding native fiber.
\end{assumption}

\begin{assumption}[Admissible quadratic regularizers]\label{ass:admissible-Q}
The regularizer \(Q\) ranges over symmetric positive-definite \(d\times d\)
matrices.  An admissible pair \((Q_1,Q_2)\) is allowed to depend only on the
researcher's chosen bridge norm or basis weights, not on unobserved potential
outcomes.
\end{assumption}

\begin{assumption}[Generic finite proxy stratum]\label{ass:generic-stratum}
The matrices \(T_P\) and \(B_P\), the vectors \(m_P\) and \(c_P\), and the
box constraints defining \(\calH_P^{\mathrm{OR}}\) and
\(\calH_P^{\mathrm{MM}}\) have fixed rank and fixed active zero/one pattern
in a neighborhood of \(P\).  Both native fibers have nonempty relative
interior in their affine hulls.
\end{assumption}

\begin{assumption}[Native finite quadratic bridge objectives]\label{ass:native-workflow}
The proximal OR bridge implementation covered here is the finite
residual-squared conditional-moment objective
\(\|T_Ph-m_P\|_\Omega^2+\alpha h^\transpose Qh\) with positive observable
weights \(\Omega\), used as a finite population abstraction of the outcome
bridge step inside proximal OR/IPW/DR analyses such as
\citet{CuiPuShiMiaoTchetgen2024JASA}.  The minimax negative-control bridge
implementation covered here is the finite-dictionary, finite-critic,
profiled quadratic specialization of \citet{KallusMaoUehara2021MinimaxNC},
with population criterion
\((B_Ph-c_P)^\transpose R_P(B_Ph-c_P)+\alpha h^\transpose Qh\) and
\(R_P\succ0\).  The two exact-fit fibers need not coincide, and no claim is
made for nonsmooth, adversarial, infinite-dimensional, or nonquadratic
minimax saddle problems that do not reduce to this finite profiled quadratic
form.
\end{assumption}

\begin{assumption}[Fixed-rank first-order regularity]\label{ass:regularity}
In a neighborhood of \(P\), the active bridge faces for
\(\calH_P^{\mathrm{OR}}\) and \(\calH_P^{\mathrm{MM}}\) are fixed, the
conditional-moment and minimax matrices are differentiable functions of the
cell probabilities, and the selected quadratic-program minimizers are unique
for each admissible \(Q\).  Under these finite-dimensional regularity
conditions, the proximal OR/IPW/DR plug-in estimator and the profiled minimax
plug-in estimator have first-order influence functions
\(\psi_{\mathrm{OR/DR},Q,P}\) and \(\psi_{\mathrm{MM},Q,P}\) obtained by
differentiating their population selected ATE maps.
\end{assumption}

\begin{assumption}[ATE-active bridge nonuniqueness]\label{ass:active-nonunique}
At least one native fiber has two elements \(h_1,h_2\in\calH_P^j\), for
\(j\in\{\mathrm{OR},\mathrm{MM}\}\), such that
\(\ell_P(h_1)\neq\ell_P(h_2)\).  Equivalently, at least one native sharp
bridge-fiber interval \(\Theta_I^j(P)\) is not a singleton.
\end{assumption}

\section{Main results}
\begin{proposition}[Support Proposition 1: null-annihilator criterion for proximal ATE point identification]\label{prop:null-annihilator}
Under Assumptions~\ref{ass:finite-bounded}--\ref{ass:fiber-sharp},
the proximal ATE for a native workflow \(j\in\{\mathrm{OR},\mathrm{MM}\}\)
is point identified if and only if
\[
\ell_P(r)=0
\quad\text{for every }r=h'-h\text{ with }h,h'\in\calH_P^j.
\]
Equivalently,
\[
\underline\theta_j=\overline\theta_j
\quad\Longleftrightarrow\quad
\ell_P \text{ is constant on } \calH_P^j.
\]
For the OR workflow this holds automatically when
\(\operatorname{ker}(T_P)=\{0\}\); for the minimax workflow it holds
automatically when \(\operatorname{ker}(B_P)=\{0\}\).  When the relevant
kernel is nonzero, point identification can still hold because the ATE
loading may annihilate all feasible native bridge-null directions.
\end{proposition}
\paragraph{Why this is new and worth studying.}
(a) This is not a novelty claim.  It is the finite linear-algebra form of a
phenomenon present in
\citet{ZhangLiMiaoTchetgen2023SPL} and
\citet{BennettKallusMaoNeweySyrgkanisUehara2023COLT}: a functional can be
identified even when the nuisance bridge is not.  The epsilon gap is that
the proposal uses this known criterion only as the support boundary inside
the native threshold \(\tau_{\mathrm{native}}(P)\); the flagship claims start
with Conjecture~\ref{conj:penalty-separation}.  (b) The concrete consumer is
\citet{CuiPuShiMiaoTchetgen2024JASA}: their proximal OR/IPW/DR estimators
can be interpreted as point-identifying estimators precisely in the invariant
OR case described here.

\begin{lemma}[Support Lemma 1: exact-fiber Tikhonov limit]\label{lem:tikhonov-limit}
Under Assumptions~\ref{ass:finite-bounded}, \ref{ass:bridge-exists},
\ref{ass:admissible-Q}, and \ref{ass:native-workflow}, for each native
workflow \(j\in\{\mathrm{OR},\mathrm{MM}\}\) and every admissible \(Q\), the
penalized selectors \(h_{\mathrm{OR},\alpha,Q}(P)\) and
\(h_{\mathrm{MM},\alpha,Q}(P)\) are the finite objectives displayed in
Section~6, and the exact-fiber selector \(h_{j,Q}(P)\) exists and is unique.
Moreover every sequence \(\alpha_n\downarrow0\) satisfies
\[
h_{j,\alpha_n,Q}(P)\longrightarrow h_{j,Q}(P),
\]
where \(h_{j,\alpha,Q}\) means \(h_{\mathrm{OR},\alpha,Q}\) when
\(j=\mathrm{OR}\) and \(h_{\mathrm{MM},\alpha,Q}\) when \(j=\mathrm{MM}\).
Hence the population regularized proximal ATE converges to
\(\theta_{j,Q}(P)=\ell_P(h_{j,Q}(P))\).
\end{lemma}
\paragraph{Why this is new and worth studying.}
(a) This is not a novelty claim.  It is the finite exact-fiber version of
the Tikhonov / penalty-selector limit familiar from
\citet{DarollesFanFlorensRenault2011Econometrica} and from weakly identified
nuisance-function analyses such as
\citet{BennettKallusMaoNeweySyrgkanisUehara2023COLT}; the epsilon gap is
that this proposal uses the limit only to name the exact finite native
population objects whose causal status is tested in
Conjectures~\ref{conj:penalty-separation}--\ref{conj:selection-range}.
(b) The concrete consumer is the finite bridge-implementation comparison
between the proximal OR/IPW/DR analysis of
\citet{CuiPuShiMiaoTchetgen2024JASA} and the finite-class profiled-quadratic
minimax specialization of \citet{KallusMaoUehara2021MinimaxNC}.

\begin{conjecture}[Conjecture 1: native equivalence frontier for finite proximal OR/DR and minimax bridge estimators (NOT YET PROVED)]\label{conj:penalty-separation}
Under Assumptions~\ref{ass:finite-bounded}--\ref{ass:generic-stratum} and
Assumptions~\ref{ass:native-workflow}--\ref{ass:regularity}, there is a
nonempty open class of finite negative-control laws \(P\) with fixed active
bridge-face pattern for which the observable threshold
\(\tau_{\mathrm{native}}(P)\) is an iff frontier for native population
equivalence:
\[
\tau_{\mathrm{native}}(P)=0
\quad\Longleftrightarrow\quad
\theta_{\mathrm{OR},Q_1}(P)=\theta_{\mathrm{MM},Q_2}(P)
\text{ for every }Q_1,Q_2\succ0 .
\]
On the left side, the common value is the singleton native causal value in
both \(\Theta_I^{\mathrm{OR}}(P)\) and \(\Theta_I^{\mathrm{MM}}(P)\), even
when a bridge itself is not unique.  Moreover, if
\(\tau_{\mathrm{native}}(P')=0\) for every law \(P'\) in the fixed-rank
neighborhood from Assumption~\ref{ass:regularity}, then the finite proximal
OR/IPW/DR plug-in estimator using \(h_{\mathrm{OR},Q_1}\) and the finite
profiled minimax plug-in estimator using \(h_{\mathrm{MM},Q_2}\) have the same
selected ATE map on that neighborhood and hence the same first-order
influence-function and efficiency-bound object for every admissible
\(Q_1,Q_2\).  On the right side,
\(\tau_{\mathrm{native}}(P)>0\) means that either at least one native
workflow has a non-singleton ATE interval, or the two native exact-fit
objectives imply different singleton ATEs; thus a single proximal point
estimate cannot be justified by native bridge invariance alone.  In that
regime, no common first-order/efficiency-bound interpretation is valid for
all finite OR/DR and profiled minimax selected ATEs unless an additional
non-native identifying restriction collapses the endpoint disagreement.  This
statement compares the two finite native objectives in
Assumption~\ref{ass:native-workflow}: the residual-squared proximal OR
conditional-moment objective motivated by
\citet{CuiPuShiMiaoTchetgen2024JASA}, and the finite-dictionary,
finite-critic, profiled-quadratic minimax specialization motivated by
\citet{KallusMaoUehara2021MinimaxNC}, before imposing any shared exact
bridge fiber.
\end{conjecture}
\paragraph{Why this is new and worth studying.}
(a) The closest papers are \citet{ZhangLiMiaoTchetgen2023SPL} on proximal
nonunique bridges and
\citet{BennettKallusMaoNeweySyrgkanisUehara2023COLT} on strongly identified
functionals, with
\citet{BennettKallusMaoNeweySyrgkanisUehara2023MinimaxIV} as the closest
regularized minimax IV selector comparator.  The flagship epsilon gap is an
observable iff frontier for the exact finite objectives in
Assumption~\ref{ass:native-workflow}: those papers show that functional
identification can survive weak or nonunique nuisance bridges, and that
penalized minimax IV can converge to a fixed solution without uniqueness or
closedness, but they do not give a finite proximal negative-control diagnostic
equivalent to common identified set, selected ATE map, influence-function
object, and efficiency-bound object on a local zero-threshold stratum for a residual-squared proximal OR/DR
conditional-moment workflow and a finite-dictionary profiled minimax workflow
in their native forms.  The
epsilon gap relative to the v2 draft is that no common exact-fiber
\(Q\)-selector is assumed; the
threshold is computed from the two native LP intervals and their endpoint
disagreement.  (b) Resolving it would give
\citet{CuiPuShiMiaoTchetgen2024JASA} users implementing the finite
residual-squared bridge objective a concrete diagnostic for when a
regularized bridge limit remains a point-identified proximal ATE despite
bridge nonuniqueness.

\begin{conjecture}[Conjecture 2: native selection range after ATE-active nonuniqueness (NOT YET PROVED)]\label{conj:selection-range}
Under Assumptions~\ref{ass:finite-bounded}--\ref{ass:active-nonunique} and
Assumption~\ref{ass:native-workflow}, on the
\(\tau_{\mathrm{native}}(P)>0\) side of
Conjecture~\ref{conj:penalty-separation} the following dichotomy is sharp.
If both native intervals are singletons, then the two finite native
objectives in Assumption~\ref{ass:native-workflow} report different
singleton ATEs under the same observed law.  Let
\[
\mathcal J(P)=\{j\in\{\mathrm{OR},\mathrm{MM}\}:
\underline\theta_j(P)<\overline\theta_j(P)\}
\]
be the set of native objectives whose sharp interval is non-singleton.  For
each \(j\in\mathcal J(P)\), there exist admissible SPD quadratic bridge
criteria \(Q_1,Q_2\) such that
\(\theta_{j,Q_1}(P)\neq\theta_{j,Q_2}(P)\).  In this penalty-dependent
regime for workflow \(j\), the native quadratic-regularizer selection range
\[
\mathcal S_j(P)=\{\theta_{j,Q}(P):Q=Q^\transpose\succ0\}
\]
has nonempty relative interior inside \(\Theta_I^j(P)\).  More sharply,
\(\overline{\mathcal S_j(P)}\) contains every ATE value \(\ell_P(h^\star)\)
of an SPD-exposed bridge point \(h^\star\in\calH_P^j\); in the
one-dimensional native bridge-fiber case,
\[
\overline{\mathcal S_j(P)}=\Theta_I^j(P).
\]
The range characterization is stable on a neighborhood \(\mathcal U\) of
\(P\) with the same finite proxy stratum and active bridge-face pattern.
Consequently a regularized proximal OR/IPW/DR bridge point value, or a
finite-class profiled-quadratic minimax bridge point value, whose finite
population bridge step has the native SPD-quadratic form can be a genuine
regularization artifact: it is a well-defined limit of the stated finite
bridge objective but not a point-identified causal estimand.
\end{conjecture}
\paragraph{Why this is new and worth studying.}
(a) The closest papers are \citet{DarollesFanFlorensRenault2011Econometrica}
and \citet{Horowitz2014ARE} on regularized inverse problems,
\citet{ZhangLiMiaoTchetgen2023SPL} on proximal nonunique bridges, and
\citet{BennettKallusMaoNeweySyrgkanisUehara2023COLT} on strongly identified
functionals.  \citet{BennettKallusMaoNeweySyrgkanisUehara2023MinimaxIV}
is a closer regularized minimax comparator because it studies convergence to a
fixed IV solution without identification or closedness; the epsilon gap is
that it does not characterize the proximal causal ATE range exposed by varying
the admissible selector.  \citet{KallusMaoUehara2021MinimaxNC} is the closest named
proximal workflow because it learns negative-control bridges under weak
uniqueness/completeness conditions, but only its finite-dictionary,
finite-critic profiled quadratic specialization is covered by the native
finite-quadratic convention.  The flagship epsilon gap is the finite-objective
equivalence frontier: none of these papers gives an observable condition
equivalent to penalty invariance versus penalty dependence for the proximal OR
residual-squared conditional-moment objective and the finite-class profiled
minimax objective in their native forms.  The
additional open-class range conjecture says that, once
\(\tau_{\mathrm{native}}(P)>0\) and ATE-active bridge nonuniqueness holds for
a native workflow, \(\mathcal S_j(P)\) has nonempty relative interior,
contains all SPD-exposed ATE values after closure, and equals the full native
sharp interval after closure when the bridge fiber is one-dimensional.  These
are genuine equivalence, coverage, and genericity claims, not the routine
fact that selected bridge values lie inside a sharp interval or that two
quadratic programs can have different minimizers.  (b) Resolving it would give
\citet{CuiPuShiMiaoTchetgen2024JASA} users, and
\citet{KallusMaoUehara2021MinimaxNC} users in the finite-class quadratic
specialization, a sharp warning: when the diagnostic fails, changing the
bridge norm can move the reported ATE through a computable part of the sharp
interval without changing the observed law or the proximal assumptions.

\begin{conjecture}[Conjecture 3: nondegenerate two-by-two negative-control witness (NOT YET PROVED)]\label{conj:binary-witness}
There exists a binary-\(W\), binary-\(Z\), no-covariate finite proximal law
satisfying Assumptions~\ref{ass:finite-bounded}--\ref{ass:generic-stratum}
with \(\Pr_P(A=a,W=w,Z=z)>0\) for all \(a,w,z\),
\(0<\E_P(Y\mid Z=z,A=a)<1\) for all \(z,a\), nonconstant conditional outcome
means across proxy-treatment cells, \(\calH_P^{\mathrm{OR}}\) a
one-dimensional line segment, and \(\ell_P\) nonconstant on that segment.
More explicitly, one candidate OR witness has
\[
T_{P,0}=
\begin{pmatrix}
1/5&4/5\\
7/10&3/10
\end{pmatrix},
\quad
m_{P,0}=\binom{3/10}{7/10},
\qquad
T_{P,1}=
\begin{pmatrix}
2/5&3/5\\
2/5&3/5
\end{pmatrix},
\quad
m_{P,1}=\binom{1/2}{1/2},
\]
where \(T_{P,a}\) maps \(h_a=(h(0,a),h(1,a))\) to
\(\E\{h(W,a)\mid Z,A=a\}\), together with the positive observable cell
masses
\[
\begin{array}{c|cc}
&W=0&W=1\\ \hline
A=0,Z=0&1/25&4/25\\
A=0,Z=1&21/100&9/100\\
A=1,Z=0&1/10&3/20\\
A=1,Z=1&1/10&3/20
\end{array}
\]
and Bernoulli conditional outcome means
\[
\E(Y\mid A=0,Z=0,W=w)=3/10,\quad
\E(Y\mid A=0,Z=1,W=w)=7/10,\quad
\E(Y\mid A=1,Z=z,W=w)=1/2.
\]
These masses imply \(\Pr_P(W=0)=9/20\) and
\(\Pr_P(W=1)=11/20\), and they reproduce the displayed \(T_{P,a}\) and
\(m_{P,a}\).  Then the
\(a=0\) bridge is unique, the \(a=1\) bridge is the bounded segment
\(\{(u,v):2u/5+3v/5=1/2,\ 0\le u,v\le1\}\), and the ATE loading
\(\ell_P(h)=\sum_w \Pr_P(W=w)\{h(w,1)-h(w,0)\}\) is nonconstant on the segment.
For this law, diagonal admissible regularizers
\[
Q_\lambda=\operatorname{diag}(\lambda,1,1,\ldots,1),\qquad \lambda>0,
\]
select a continuous nonconstant curve
\(\lambda\mapsto\theta_{\mathrm{OR},Q_\lambda}(P)\) whose image has positive
length inside \(\Theta_I^{\mathrm{OR}}(P)\).  In the displayed \(a=1\) fiber, writing
\(u=h(0,1)\) gives \(v=(5-4u)/6\) and \(u\in[0,5/4]\cap[0,1]=[0,1]\);
the diagonal \(Q_\lambda\) selector solves
\[
\min_{0\le u\le1}\lambda u^2+\left(\frac{5-4u}{6}\right)^2,
\]
with \(u_\lambda=1\), \(v_\lambda=1/6\) for \(0<\lambda\le 1/9\), and
\[
u_\lambda=\frac{5}{4+9\lambda},\qquad
v_\lambda=\frac{15\lambda}{2(4+9\lambda)}
\]
for \(\lambda\ge1/9\).
Since \(\Pr_P(W=0)=9/20\), the selected treated bridge mean is
\[
\frac{9}{20}u_\lambda+\frac{11}{20}v_\lambda
\]
which is nonconstant in \(\lambda\); the untreated bridge mean is fixed by
the unique \(a=0\) solution.
\end{conjecture}
\paragraph{Why this is new and worth studying.}
(a) \citet{MiaoGengTchetgen2018Biometrika} and
\citet{ShiMiaoNelsonTchetgen2020JRSSB} use finite or categorical proxy
structures to motivate point identification under rank restrictions, but do
not give a minimal finite witness where bridge regularization changes the
causal ATE.  The epsilon gap is a small auditable negative-control example
showing that penalty dependence is not merely an infinite-dimensional
pathology.  (b) The concrete consumer is
\citet{GhassamiShpitserTchetgen2023ProxyPI}: a binary witness would let
proxy partial-ID bounds and regularized proximal point estimates be compared
on the same observed law.

\section{Sanity-check corollaries}
\begin{corollary}[Completeness corner]\label{cor:completeness}
Under Assumptions~\ref{ass:finite-bounded}--\ref{ass:fiber-sharp}, if
\(\operatorname{ker}(T_P)=\{0\}\), then
\(\calH_P^{\mathrm{OR}}\) is a singleton,
\(\Theta_I^{\mathrm{OR}}(P)\) is a singleton, and
\(\theta_{\mathrm{OR},Q}(P)\) is the same for every admissible \(Q\).  This
is the finite-support reduction to the Miao--Geng--Tchetgen
point-identification regime for the proximal OR workflow.
\end{corollary}

\begin{corollary}[Strong functional identification with nonunique bridges]\label{cor:strong-functional}
Under Assumptions~\ref{ass:finite-bounded}--\ref{ass:fiber-sharp}, if there
exists a vector \(\varphi\) such that \(\ell_P(h)=\varphi^\transpose T_Ph\)
for all \(h\) in the affine hull of \(\calH_P^{\mathrm{OR}}\), then
\(\Theta_I^{\mathrm{OR}}(P)\) is a singleton even if
\(\operatorname{ker}(T_P)\neq\{0\}\).  This is the finite adjoint-range
reduction to the Santos and Bennett et al. strong-functional identification
condition.
\end{corollary}

\begin{corollary}[Penalty-artifact corner]\label{cor:artifact}
Under Assumptions~\ref{ass:finite-bounded}--\ref{ass:active-nonunique}, if
Conjecture~\ref{conj:selection-range} holds for \(P\), then there exist
either two native objectives with different singleton ATEs or two admissible
regularizers \(Q_1,Q_2\) inside one native workflow \(j\) with
\(\theta_{j,Q_1}(P)\neq\theta_{j,Q_2}(P)\).  Hence at least one selected
value differs from any claimed single proximal ATE value.  The reduction is
immediate because the same observed law and the same proximal assumptions
generate two selected values.
\end{corollary}

\begin{corollary}[Singleton-disagreement boundary]\label{cor:singleton-disagreement}
Under Assumptions~\ref{ass:finite-bounded}--\ref{ass:native-workflow}, if
\(\Theta_I^{\mathrm{OR}}(P)=\{\theta_{\mathrm{OR}}^\star\}\),
\(\Theta_I^{\mathrm{MM}}(P)=\{\theta_{\mathrm{MM}}^\star\}\), and
\(\theta_{\mathrm{OR}}^\star\neq\theta_{\mathrm{MM}}^\star\), then
\[
\tau_{\mathrm{native}}(P)=|\theta_{\mathrm{OR}}^\star-\theta_{\mathrm{MM}}^\star|>0.
\]
The LP certificate is the pair of singleton native bridge intervals: each
workflow is internally invariant over \(Q\), but the endpoint-disagreement
terms in \(\tau_{\mathrm{native}}\) certify cross-objective non-equivalence.
\end{corollary}

\section{Proof-sketch route}
Proposition~\ref{prop:null-annihilator} reduces to affine linear algebra:
write each native fiber as \(h_0+\operatorname{ker}(A_P^j)\) intersected
with box constraints and check whether the linear functional \(\ell_P\) is
constant on that convex polytope.
Lemma~\ref{lem:tikhonov-limit} follows from compactness of \([0,1]^d\),
strict convexity of \(h^\transpose Qh\), and the standard exact-penalty limit
for minimizing either \(\|T_Ph-m_P\|^2+\alpha h^\transpose Qh\) or
\((B_Ph-c_P)^\transpose R_P(B_Ph-c_P)+\alpha h^\transpose Qh\) as
\(\alpha\downarrow0\); it is not a novelty target and merely connects those
standard limits to each native exact bridge fiber.
The flagship conjectures should be attacked by first solving the four
observable LPs for
\((\underline\theta_{\mathrm{OR}},\overline\theta_{\mathrm{OR}},
\underline\theta_{\mathrm{MM}},\overline\theta_{\mathrm{MM}})\), proving the
linear-algebra equivalence that each native LP interval is a singleton iff
its \(\kappa_j(P)\) vanishes, then proving that
\(\tau_{\mathrm{native}}(P)=0\) is exactly the additional cross-workflow
agreement condition.  Under the fixed-rank regularity assumption, differentiate
the finite QP solution maps by the KKT system to compare
\(\psi_{\mathrm{OR/DR},Q,P}\) and \(\psi_{\mathrm{MM},Q,P}\); the conjectured
flagship boundary says these first-order objects and their efficiency-bound
interpretations coincide exactly when the zero-threshold equality persists on
the fixed-rank stratum and fail to admit a common native interpretation on the
positive-threshold side.  The
non-invariant side adds Assumption~\ref{ass:active-nonunique} and
characterizes which exposed bridge points satisfy the SPD first-order selector
condition in each native fiber by the finite semidefinite feasibility
certificate \(-2Qh^\star\in N_{\calH_P^j}(h^\star)\), \(Q\succ0\).
The final range step projects those exposed points through \(\ell_P\) and
uses fixed-rank active-set stability to promote the selection range to an
open class of observable laws.  The binary witness
requires the displayed two-proxy matrix calculation: solve the unique
\(a=0\) bridge, enumerate the \(a=1\) one-dimensional bridge fiber, compute
the diagonal \(Q_\lambda\) minimizer in closed form, and verify that the
resulting ATE loading changes with \(\lambda\) while the same finite proximal
restrictions remain satisfied.

\section{Honest scope flags}
Proposition~\ref{prop:null-annihilator} and Lemma~\ref{lem:tikhonov-limit} are
routine finite-dimensional statements; the lemma is standard support from
inverse-problem regularization and weak-nuisance work, not a novelty claim.
The real novelty is split across Conjecture~\ref{conj:penalty-separation}
and Conjecture~\ref{conj:selection-range}.  The first is the finite
native-equivalence frontier for selected ATE maps, identified sets, and
first-order/efficiency-bound objects; it deliberately does not assume
Assumption~\ref{ass:active-nonunique}; the second handles the non-invariant
native selection range after ATE-active nonuniqueness is imposed.
Conjecture~\ref{conj:binary-witness} is also open; it is meant to force a
small auditable example before any broader theorem is claimed.  This pivot
withdraws the rejected angle's stronger claim that a scalar smallest singular
value alone creates sharp point/interval/Manski phase boundaries.  The draft
also assumes finite bridge-fiber sharpness rather than proving it, because
the prior \texttt{pid\_mediation\_proxy} bank shows that sharpness is a
separate realization problem.  The v4 revision narrows every
\citet{CuiPuShiMiaoTchetgen2024JASA} and
\citet{KallusMaoUehara2021MinimaxNC} comparison to the finite native
objectives explicitly displayed here; it no longer claims that the full
published estimators are identical to those finite quadratic limits.

\section{Iteration trail}
\begin{itemize}[leftmargin=*]
\item v1 (pivot) -- initial draft for the regularization-artifact angle after
the finite sharp-set phase-modulus angle was archived; prior verdict: none.
\item v2 (revise) -- demoted the exact-fiber Tikhonov limit to standard
support, split the flagship frontier into invariant and non-invariant
conjectures, and added the binary-witness selector calculation; prior verdict:
REVISE.
\item v3 (revise) -- moved the null-annihilator criterion out of the novelty
payload, replaced the shared exact-fiber assumption with native OR and
minimax objectives, and recast the flagship threshold as
\(\tau_{\mathrm{native}}(P)\); prior verdict: REVISE.
\item v4 (revise) -- pinned the finite residual-squared OR and profiled
quadratic minimax objectives, fixed the Support Lemma 1 structure and
minimax selector notation, and made Conjecture 2 workflow-specific when only
one native interval is non-singleton; prior verdict: REVISE.
\item v5 (revise) -- added the missed minimax-IV comparator, upgraded
Conjecture 1 to a finite OR/DR-vs-minimax native-equivalence frontier for
identified sets and first-order/efficiency-bound objects, and made the
LP/QP/SDP computation recipe explicit; prior verdict: REVISE.
\end{itemize}

\section*{Bibliography}
\begin{thebibliography}{99}
\bibitem[Bennett et~al.(2023a)]{BennettKallusMaoNeweySyrgkanisUehara2023COLT}
Bennett, A., Kallus, N., Mao, X., Newey, W., Syrgkanis, V., and Uehara, M.
(2023).
\newblock Inference on strongly identified functionals of weakly identified
functions.
\newblock \emph{Proceedings of Machine Learning Research}, 195:2265--2265.

\bibitem[Bennett et~al.(2023b)]{BennettKallusMaoNeweySyrgkanisUehara2023MinimaxIV}
Bennett, A., Kallus, N., Mao, X., Newey, W., Syrgkanis, V., and Uehara, M.
(2023).
\newblock Minimax instrumental variable regression and \(L_2\) convergence
guarantees without identification or closedness.
\newblock \emph{Proceedings of Machine Learning Research}, 195:2291--2318.

\bibitem[Canay et~al.(2013)]{CanaySantosShaikh2013Econometrica}
Canay, I. A., Santos, A., and Shaikh, A. M. (2013).
\newblock On the testability of identification in some nonparametric models
with endogeneity.
\newblock \emph{Econometrica}, 81(6):2535--2559.

\bibitem[Cui et~al.(2024)]{CuiPuShiMiaoTchetgen2024JASA}
Cui, Y., Pu, H., Shi, X., Miao, W., and Tchetgen Tchetgen, E. J. (2024).
\newblock Semiparametric proximal causal inference.
\newblock \emph{Journal of the American Statistical Association},
119(546):1348--1359.

\bibitem[Darolles et~al.(2011)]{DarollesFanFlorensRenault2011Econometrica}
Darolles, S., Fan, Y., Florens, J.-P., and Renault, E. (2011).
\newblock Nonparametric instrumental regression.
\newblock \emph{Econometrica}, 79(5):1541--1565.

\bibitem[Ghassami et~al.(2023)]{GhassamiShpitserTchetgen2023ProxyPI}
Ghassami, A., Shpitser, I., and Tchetgen Tchetgen, E. J. (2023).
\newblock Partial identification of causal effects using proxy variables.
\newblock \emph{arXiv:2304.04374}.

\bibitem[Horowitz(2014)]{Horowitz2014ARE}
Horowitz, J. L. (2014).
\newblock Ill-posed inverse problems in economics.
\newblock \emph{Annual Review of Economics}, 6:21--51.

\bibitem[Kallus et~al.(2021)]{KallusMaoUehara2021MinimaxNC}
Kallus, N., Mao, X., and Uehara, M. (2021).
\newblock Causal inference under unmeasured confounding with negative
controls: A minimax learning approach.
\newblock \emph{arXiv:2103.14029}.

\bibitem[Manski(1990)]{Manski1990AER}
Manski, C. F. (1990).
\newblock Nonparametric bounds on treatment effects.
\newblock \emph{American Economic Review}, 80(2):319--323.

\bibitem[Miao et~al.(2018)]{MiaoGengTchetgen2018Biometrika}
Miao, W., Geng, Z., and Tchetgen Tchetgen, E. J. (2018).
\newblock Identifying causal effects with proxy variables of an unmeasured
confounder.
\newblock \emph{Biometrika}, 105(4):987--993.

\bibitem[Santos(2011)]{Santos2011JoE}
Santos, A. (2011).
\newblock Instrumental variable methods for recovering continuous linear
functionals.
\newblock \emph{Journal of Econometrics}, 161(2):129--146.

\bibitem[Santos(2012)]{Santos2012Econometrica}
Santos, A. (2012).
\newblock Inference in nonparametric instrumental variables with partial
identification.
\newblock \emph{Econometrica}, 80(1):213--275.

\bibitem[Shi et~al.(2020)]{ShiMiaoNelsonTchetgen2020JRSSB}
Shi, X., Miao, W., Nelson, J. C., and Tchetgen Tchetgen, E. J. (2020).
\newblock Multiply robust causal inference with double-negative control
adjustment for categorical unmeasured confounding.
\newblock \emph{Journal of the Royal Statistical Society: Series B},
82(2):521--540.

\bibitem[Tchetgen Tchetgen et~al.(2024)]{TchetgenYingCuiShiMiao2024StatSci}
Tchetgen Tchetgen, E. J., Ying, A., Cui, Y., Shi, X., and Miao, W. (2024).
\newblock An introduction to proximal causal inference.
\newblock \emph{Statistical Science}, 39(3):375--390.

\bibitem[Zhang et~al.(2023)]{ZhangLiMiaoTchetgen2023SPL}
Zhang, J., Li, W., Miao, W., and Tchetgen Tchetgen, E. J. (2023).
\newblock Proximal causal inference without uniqueness assumptions.
\newblock \emph{Statistics \& Probability Letters}, 198:109836.
\end{thebibliography}

\end{document}
